\section{Benchmark}

\subsection{Scenes}
The corpus has 900 physical scenes rendered in CARLA Town10HD. Scenes fall
into 18 cells from a full factorial of three factors: viewpoint disparity
(the angular gap between the two cameras on a shared orbit), the presence of a
static occluder, and illumination (ClearNoon, CloudySunset, WetNoon). Each
cell holds 50 scenes. The generator is seeded (seed 42) and its manifest is
hashed, so a run can be replayed. Each scene is rendered once and read by
every detector pairing.

Two agents view each scene from cameras 14\,m above ground (800$\times$600,
$90^\circ$ field of view). The scene supplies the agents with the same objects
and the same static geometry, which is the shared context we want to measure.

\subsection{Detectors and pairings}
Each of the 1{,}800 frames (900 scenes $\times$ two viewpoints) goes through
both YOLO11m and RT-DETR-L, giving 3{,}600 detector reports. Detections are
matched to truth at IoU 0.5. Three pairings are scored: YOLO11m with
YOLO11m, RT-DETR-L with RT-DETR-L, and YOLO11m with RT-DETR-L. A task is one
scene under one pairing, so the corpus has $900\times 3 = 2{,}700$ tasks.

\subsection{Truth and visibility}
Ground truth comes from CARLA's instance-segmentation pixels, not from
projected 3D boxes. An object enters an agent's visible set only if it covers
enough pixels in that agent's frame. We score only objects visible to both
agents. An object one agent cannot see says nothing about joint failure, and
including it would pull the double-fault rate down. Projected boxes cannot
tell an occluded object from a visible one, which is why we avoided them.
Truth is kept in a separate directory and never reaches the policies.
